\begin{remark}[Two hardening paths and \texorpdfstring{$(u,q,\varepsilon)$}{(u,q,eps)} scheduling]\label{rem:pom-hardening-u-vs-q}
The ``hardening'' of order gates admits at least two executable paths within the interface of this paper:
\begin{itemize}
  \item \textbf{Temperature path:} Lowering \(u\) (Remark \ref{rem:pom-order-as-zero-temp}) causes the soft order gate \(\mathrm{PROJ}_u\) to harden to \(P_{\le}\); the time-side cost is primarily reflected in the twisted spectral gap/mixing scale \(\tau_{\mathrm{mix}}(u)\) (Proposition \ref{prop:pom-rewrite-audit-loop}).
  \item \textbf{Moment-quantile path:} Fixing (or not drastically lowering) \(u\), one raises \(q\) and uses the moment-kernel (\(\mathrm{MOM}_q\)) to produce a tail certificate, thereby triggering \((\mathrm{RCRT}_\varepsilon)\) (Proposition \ref{prop:pom-rcrt-epsilon}) to compile the order gate into a spatialized residue+CRT readout; the space-side cost is primarily reflected in the prime-register quota \(\log P\) (Corollary \ref{cor:pom-prime-register-quantile-budget} and the \(\gamma\)-certificate of Remark \ref{rem:pom-quantile-budget-computable}).
\end{itemize}
Thus the compiler can treat \((u,q,\varepsilon)\) as three knobs: \(u\) primarily controls the time scale, \(q\) primarily controls tail amplification and spatial quota, and \(\varepsilon\) primarily controls the risk budget; the additive law of Proposition \ref{prop:pom-eps-sound-compose} is used to account for the overall failure probability across multiple approximate order-gate eliminations.
The two-dimensional pressure surface \(P(q,u)\) provides a unified coordinate system placing both paths into a single ``spectral phase diagram'' (\S\ref{subsec:pom-pressure-surface-qu}).
\end{remark}

\begin{table}[H]
\centering
\scriptsize
\setlength{\tabcolsep}{6pt}
\caption{Space--time duality diagram: closed-form budgets for the same prime-register capacity variable across two readout tasks. The first row gives the spatial quota for ``unique fiber-label reconstruction'' (controlled by the collision moment spectrum \(\{r_q\}\)); the second row gives the temporal budget for ``cylinder congruence equidistribution'' (controlled by the worst-case twisted spectral gap \(\Lambda_m/\lambda_m\) of the profinite Chebotarev mechanism; see Corollary \ref{cor:pom-profinite-cylinder-effective-n} and the inverse Corollary \ref{cor:pom-profinite-cylinder-capacity}).}
\label{tab:pom-space-time-duality-prime-register}
\begin{tabular}{p{0.25\linewidth} p{0.60\linewidth} p{0.12\linewidth}}
\toprule
Readout task & Sufficient budget (closed-form inequality family) & Fingerprint input\\
\midrule
Quantile-CRT (fiber-label exceedance failure \(\le\varepsilon\)) &
\(
P_m>2B_m^{(q)}(\varepsilon),
\ \
B_m^{(q)}(\varepsilon)=\varepsilon^{-1/(q-1)}(r_q/2)^{m/(q-1)}
\)
(Corollary \ref{cor:pom-prime-register-q-quota-family}; finite-scale version uses \(S_q(m)\) to obtain \(B_{m,q}^{\mathrm{fin}}\)) &
\(\{r_q\}\)\\

Chebotarev cylinder (relative error \(\le\delta\)) &
If \(\Lambda_m\le \lambda_m^{1/2-\varepsilon_{\mathrm{gap}}}\), then
\(
n\gtrsim \log(|G_m|/\delta)/((1/2+\varepsilon_{\mathrm{gap}})\log\lambda_m)
\)
(Corollary \ref{cor:pom-profinite-cylinder-effective-n}), equivalently
\(
\log|G_m|\lesssim (1/2+\varepsilon_{\mathrm{gap}})n\log\lambda_m-\log(1/\delta)
\)
(Corollary \ref{cor:pom-profinite-cylinder-capacity}) &
\(\Lambda_m/\lambda_m\)\\
\bottomrule
\end{tabular}
\end{table}

\noindent\textbf{Splicing into a single ``space--time strip''}
On the prime-register axis, take \(G_m=\ZZ/P_m\ZZ\) (so that \(|G_m|=P_m\)).
Then the two budget rows of Table \ref{tab:pom-space-time-duality-prime-register} splice directly through the same capacity variable \(\log P_m\): the space side yields
\(
\log P_m\gtrsim \log 2+m\beta_q+\frac{1}{q-1}\log(1/\varepsilon)
\)
(Corollary \ref{cor:pom-prime-register-q-quota-family}), while the time side yields
\(
\log P_m\lesssim (1/2+\varepsilon_{\mathrm{gap}})n\log\lambda_m-\log(1/\delta)
\);
together they form an auditable feasibility strip in the \((n,\log P_m)\) plane: the same register is simultaneously constrained by the two spectral fingerprint families \(\{r_q\}\) and \(\Lambda_m/\lambda_m\).
\par
If the cylinder's empirical readout is realized via fold histograms, then the time side incurs, in addition to the dynamical tail term, a finite-sample term; under the rotational scan (Kronecker/Sturmian) model, this sample term can be deterministically controlled by the discrepancy \(D_N^*\), thereby yielding a fully deterministic \((m,n,N)\) three-budget law (Corollary \ref{cor:auditable-profinite-chebotarev-three-budget-deterministic}).

\begin{remark}[Spectral-tax calibration and ``adding axes does not necessarily accelerate'']\label{rem:pom-spectral-tax-workload}
The explicit relative error bound of Corollary \ref{cor:pom-profinite-cylinder-effective-n} in the worst case (\(|A|=1\)) can be read as
\[
\mathrm{Err}_m(n)\ \lesssim\ |G_m|\left(\frac{\Lambda_m}{\lambda_m}\right)^n.
\]
Let
\(
r_m:=\Lambda_m/\lambda_m\in(0,1]
\)
and the spectral-tax rate
\(
\Delta_m:=-\log r_m
\)
(i.e., the step-count mixing scale \(\tau_{\mathrm{mix}}(m):=1/\Delta_m\)).
Then provided that
\[
n\ \gtrsim\ \frac{\log|G_m|+\log(1/\delta)}{\Delta_m},
\]
the relative error can be suppressed to \(O(\delta)\) (absorbing constants into \(\gtrsim\)).
Consequently, ``using more congruence/character axes (space) to replace longer primitive lengths (time)'' is not free: each additional layer simultaneously changes
\(|G_m|\) and \(r_m\), and the net effect can be directly audited via the workload fingerprint
\[
W(m):=\frac{\log|G_m|}{\Delta_m}=\log|G_m|\cdot\tau_{\mathrm{mix}}(m)
\]
: if adding an axis increases \(|G_m|\) but drives \(r_m\to 1\) (spectral gap shrinks, \(\Delta_m\downarrow\)), then \(W(m)\) may increase, resulting in ``slower'' performance.
The scan table for the real-input \(40\)-state kernel on discrete prime modulus axes (Table \ref{tab:real-input-40-spectral-tax-workload}) provides a reproducible direct illustration (e.g., the comparison between \(p=3\) and \(p=5\)).
\end{remark}

\input{sections/generated/tab_real_input_40_spectral_tax_workload}

\begin{remark}[Maximum modulus control law: multiplicative space, quadratic time (near-coboundary + falsifiable selection law)]\label{rem:pom-max-modulus-control-law}
This paragraph further splices the ``multiplicative space budget'' and ``time spectral-tax'' of Table \ref{tab:pom-space-time-duality-prime-register} into a single testable engineering law: on a separable direct-product residue axis, the \emph{spatial capacity} of the prime-register is determined by the product \(P=|G|\), while the \emph{temporal diffusion} at leading order is determined by the ``worst small angle/maximum modulus,'' not by \(P\) itself.
\par
\medskip\noindent\textbf{Setup (direct-product residue axis)}
Take pairwise coprime moduli \(\{p_1,\dots,p_k\}\) and let
\[
G:=\prod_{i=1}^k\ZZ/p_i\ZZ,
\qquad
P:=|G|=\prod_{i=1}^k p_i,
\qquad
L:=\max_i p_i.
\]
Denote the worst-case twisted ratio for ``cylinder congruence equidistribution'' by
\(
r_G:=\max_{\chi\neq\chi_0}\rho(M_\chi)/\lambda
\)
and let the spectral-tax rate be
\(
\Delta_G:=-\log r_G
\)
(consistent with the notation of Remark \ref{rem:pom-spectral-tax-workload}).
\par
\medskip\noindent\textbf{(Small-angle quadratic law)}
In the near-coboundary small-angle regime, the twisted gap is controlled by the variance density: for any one-dimensional direction \(F\) with minimal nontrivial angle \(2\pi/m\),
\[
\log\frac{\rho(M_{it_mF})}{\lambda}
=-\frac{\sigma_F^2}{2}\left(\frac{2\pi}{m}\right)^2+O(m^{-4})
\]
(Corollary \ref{cor:dirichlet-modulus-diffusion-barrier}; where \(\sigma_F^2=a^{\mathsf T}\Sigma a\), \(\Sigma=\nabla^2P(0)\)).
Thus, the leading-order spectral gap for the ``pure-axis twist'' along the \(i\)-th axis (nontrivial only on \(\ZZ/p_i\), with \(j=\pm1\)) is
\[
\Delta_i\ \approx\ 2\pi^2\,\sigma_i^2/p_i^2
=4\pi^2\,\kappa_i/p_i^2,
\qquad \kappa_i:=\sigma_i^2/2.
\]
\par
\medskip\noindent\textbf{(Falsifiable selection law \(\Rightarrow\) maximum modulus control)}
If the direct-product axis satisfies the ``single-axis slow-mode selection law'':
\[
\arg\max_{\chi\neq\chi_0}\frac{\rho(M_\chi)}{\lambda}
\ \text{is attained by a pure-axis character (only one axis nontrivial), and occurs at the largest modulus }p_{i^\star}=L,
\]
then the leading order of \(\Delta_G\) is locked to this worst small angle, and consequently
\[
n\ \gtrsim\ \frac{\log|G|+\log(1/\delta)}{\Delta_G}
\ \approx\
\frac{L^2}{4\pi^2\,\kappa_{i^\star}}\Bigl(\log P+\log\tfrac{1}{\delta}\Bigr).
\]
This shows that, under the above auditable premise, \textbf{the spatial budget of the prime-register is multiplicative \(P\), but the leading-order time budget is quadratic \(L^2\), entering only through \(\log P\)}.
\par
The ``single-axis slow-mode selection law'' itself is falsifiable/auditable: for instance, in the \(((3,3,p))\) pure-collision-axis scan of the real-input 40-state kernel, the worst character index is stably \((0,0,\pm1)\) (Table \ref{tab:real_input_40_arity_335_n2_selection_law_primes}), consistent with the near-coboundary small-angle mechanism; the more general selection law under the quadratic approximation is equivalent to a metric shortest vector problem (Theorem \ref{thm:real-input-40-near-coboundary-svp}).
\end{remark}

\begin{remark}[Primorial modulus selection and \texorpdfstring{$\mathrm{poly}(\log P)$}{poly(log P)} time closure (external number-theoretic fact)]\label{rem:pom-primorial-polylog-time}
Under the spatial constraint given by Corollary \ref{cor:pom-prime-register-quantile-budget},
\(
P>2B_m(\varepsilon)=2\exp(m\gamma_m(\varepsilon))
\),
the capacity's leading scale satisfies \(\log P\asymp m\gamma_m(\varepsilon)\).
\par
If one wishes to minimize the largest modulus \(L=\max_i p_i\) while satisfying \(\prod_i p_i\ge P\), number theory (the linear-order growth of the Chebyshev \(\vartheta\) function) shows that for a given cutoff \(L\), the product \(\prod_{p\le L}p\) is maximized; moreover
\[
\log\!\left(\prod_{p\le L}p\right)=\vartheta(L)=\Theta(L).
\]
Therefore, when adopting the primorial scheme of ``taking all primes \(\le L\)'', one has
\[
L\ =\ \Theta(\log P)\ =\ \Theta(m\gamma_m(\varepsilon)).
\]
Substituting this back into the time budget of Remark \ref{rem:pom-max-modulus-control-law} (and treating \(\kappa_{i^\star}\) as a constant extractable via the spectral scan/Poisson equation interface; see Remark \ref{rem:arity-335-kappa-poisson}), one obtains the leading-order complexity law
\[
n\ \gtrsim\ \frac{(\log P)^2}{4\pi^2\kappa_{i^\star}}\Bigl(\log P+\log\tfrac{1}{\delta}\Bigr).
\]
In particular, when \(\delta\) is a constant-level (or polynomial-level) target precision, the time cost exhibits an explicit \(\,O\big((\log P)^3\big)\,\) scaling: \textbf{exponential in space (Euler product), but only polynomial in \(\log P\) for time}.
\end{remark}

\begin{remark}[Variance-matched modulus allocation]
Regarding the register design law derived from the near-coboundary quadratic law (i.e., the modulus allocation criterion that optimizes the diffusion-time upper bound using the pressure Hessian/covariance density $\Sigma$), its detailed definition and the minimax impedance matching principle are given in Remark \ref{rem:pom-variance-matched-moduli}.
\end{remark}

\begin{remark}[Three-term decomposition of the budget exponent: geometric inevitable loss + R\'enyi non-uniformity penalty + confidence term]\label{rem:pom-quantile-budget-three-term}
The Chernoff certificates of this subsection can be viewed both as an explicit transform ``spectral/moment side \(\Rightarrow\) prime-register spatial budget'' and can be further decomposed into three \emph{computable} terms, thereby quantifying the ``space/time'' trade-off as an additive law.

\medskip\noindent\textbf{(i) Exact finite-scale identity (no asymptotic assumption required)}
Let \(u_m\) be the uniform distribution on \(X_m\): \(u_m(x)=1/|X_m|\).
For any \(q\ge 2\), denote by \(\mathrm{D}_q(w_m\Vert u_m)\) the R\'enyi-\(q\) divergence:
\[
\mathrm{D}_q(w_m\Vert u_m)
:=\frac{1}{q-1}\log\Big(\sum_{x\in X_m}w_m(x)^q\cdot |X_m|^{q-1}\Big)
\ \ge\ 0.
\]
Then for any \(\varepsilon\in(0,1)\) and any \(q\ge 2\), the finite-scale certificate of Proposition \ref{prop:pom-quantile-chernoff}
\[
\gamma^{\mathrm{cert,fin}}_{m,q}(\varepsilon)
:=\frac{\frac{1}{m}\log S_q(m)-\log 2+\frac{1}{m}\log(1/\varepsilon)}{q-1}
\]
satisfies the following \textbf{exact decomposition}:
\[
\boxed{\;
\gamma^{\mathrm{cert,fin}}_{m,q}(\varepsilon)
=
\underbrace{\frac{1}{m}\log\Big(\frac{2^m}{|X_m|}\Big)}_{\text{geometric inevitable loss (mean fiber exponent)}}
\;+\;
\underbrace{\frac{1}{m}\mathrm{D}_q(w_m\Vert u_m)}_{\text{non-uniformity penalty (output distribution deviation from uniform)}}
\;+\;
\underbrace{\frac{1}{q-1}\cdot\frac{1}{m}\log(1/\varepsilon)}_{\text{confidence term (tail reliability cost)}}
\;}
\]
The proof is a purely algebraic identity: substitute \(\sum_x w_m(x)^q=S_q(m)/2^{qm}\) into \(\mathrm{D}_q\) and simplify.

\medskip\noindent\textbf{(ii) Asymptotic three-term decomposition (aligned with the \(\{r_q\}\) bridge-kit)}
When \(S_q(m)=\Theta(r_q^m)\) and \(|X_m|=F_{m+2}\sim\varphi^m\), the above expression converges at the exponential scale to
\[
\gamma(q,\alpha)
=
\log\frac{2}{\varphi}
\;+\;
\Delta_q
\;+\;
\frac{\alpha}{q-1},
\qquad
\Delta_q:=\lim_{m\to\infty}\frac{1}{m}\mathrm{D}_q(w_m\Vert u_m)
:=\frac{\log r_q-q\log 2+(q-1)\log\varphi}{q-1},
\]
where \(\alpha:=\lim_{m\to\infty}\frac{1}{m}\log(1/\varepsilon_m)\) represents the exponential-level reliability requirement \(\varepsilon_m\approx e^{-\alpha m}\).
In particular, for \(q=2\),
\(
\Delta_2=\log(\varphi r_2/4)
\)
is precisely the second-order deviation exponent constant appearing throughout this paper (cf.\ the constant-term calibration of Corollary \ref{cor:pom-shannon-entropy-squeeze}).
\end{remark}

\subsubsection{Implementable interfaces: minimal pipeline for anomaly signatures and moment certificates}\label{subsec:pom-implementable-interfaces}

To bring ``rewrite certificates/anomaly signatures/spectral identity tests'' to a reproducible execution layer, this paragraph provides a minimal pipeline in one-to-one correspondence with the repository scripts.
The core principle is: all signatures in this paper originate from a single closed loop
$$
\text{finite kernel/matrix family}\ \Longrightarrow\ \zeta(z)=\det(I-zM)^{-1}\ \Longrightarrow\ (\text{twist})\ \Longrightarrow\ \log\mathfrak{M}\ \Longrightarrow\ \text{signature}.
$$

\paragraph{Anomaly signature interface: \texorpdfstring{$\mathsf{Anom}_G$}{Anom} (twisted finite-part anomaly signature)}
For Definition \ref{def:pom-anom-signature}, engineering-wise only two steps are needed: enumerate nontrivial characters and compute the corresponding \(\log\mathfrak{M}_\chi(\theta)\).
This paper already has a script implementing the one-parameter family \(\theta=(t,0,-t)\) on the \(\chi\)-slice (i.e., injecting arity-charge as potential) and outputting \((P_\chi(t),\log\mathfrak{M}_\chi(t))\):
\texttt{scripts/exp\_sync\_kernel\_real\_input\_40\_logM\_chi.py}.
Its internal implementation directly uses the closed form of Proposition \ref{prop:abel-finite-part-mobius-zeta} (M\"obius tail truncation) and computes residue constants from matrix eigenvalues (see the script header comments).

\begin{remark}[Minimal executable workflow for anomaly signatures (pseudocode)]\label{rem:pom-anom-pseudocode}
Given kernel \(K\), lift group \(G\), potential parameter \(\theta\) (or temperature slice \(t\)), output \(\mathsf{Anom}_G(K;\theta)\):
\begin{verbatim}
Input: kernel K (weighted matrix builder M_theta), finite abelian label group G,
       parameter theta, mobius tail cutoff k_max
Output: Anom := { logM_chi(theta) : chi in G^ \ {chi0} }

for chi in characters(G) excluding trivial:
    build twisted weighted matrix M_theta,chi  (character-twist on edges/labels)
    lambda := spectral_radius(M_theta,chi)
    z_star := 1/lambda
    logC := log residue constant at z_star (eigenvalue product formula)
    logM_chi := logC + sum_{k=2..k_max} mu(k)/k * log zeta_theta,chi(z_star^k)
return Anom
\end{verbatim}
Here \(\log\zeta=-\log\det(I-zM)\) can be computed in a numerically stable manner using \texttt{slogdet} (as in the script implementation).
If only ``identity testing/zero testing'' is needed, one may compare the \(\log\mathfrak{M}_\chi(\theta)\) vectors at a small number of sample points \(\theta\in\Theta_{\mathrm{grid}}\) (Corollary \ref{cor:pom-three-gen-normal-form-decidable}).
\end{remark}

\paragraph{Slow-mode certificate interface: twisted spectral radius ratio and worst character (slow-mode certificate)}
The anomaly signature is a constant-level residual; exponential-level slow-mode certificates are given by \(\rho_\chi/\lambda\).
For certain potentials of the real-input 40-state kernel (e.g., output potential), this paper has already provided full tables of ``root-of-unity twists'' and worst-character outputs:
\texttt{scripts/exp\_real\_input\_40\_output\_potential\_dirichlet\_twists.py},
with output table \texttt{sections/generated/tab\_real\_input\_40\_output\_potential\_dirichlet\_twists.tex}
reporting the worst \(j\) and \(\rho_m/\lambda\) for each modulus \(m\).
This interface can be used alongside \(\mathsf{Anom}\): the former gives the exponential gap, the latter gives the constant drift.

\paragraph{Tail budget interface: moment recursion certificates (\texorpdfstring{$r_q$}{rq} and tail budgets)}
For the ``auditable certificates'' of the collision moments \(S_q(m)\), this paper employs two equivalent calibrations:
\begin{itemize}
  \item \textbf{Recursion certificate:} Verify \(S_q\) using integer recursion coefficients (Table \ref{tab:fold_collision_moment_recursions}).
  The corresponding script is \texttt{scripts/exp\_fold\_collision\_moment\_recursions.py} (enumerates \(d_m\), verifies the recursion, and outputs JSON and LaTeX tables).
  \item \textbf{Spectral certificate:} Obtain the Perron root \(r_q\) from the companion matrix or moment-kernel realization of the recursion, and derive pressure/tail upper bounds accordingly (Proposition \ref{prop:pom-quantile-chernoff} and Remark \ref{rem:pom-quantile-budget-computable}).
  \item \textbf{Finite-part/moment-anomaly certificate:} Reuse the Abel/Hadamard finite-part closed form (M\"obius tail truncation) of Definition \ref{def:abel-finite-part} on the \(\zeta\)-interface of the \(q\)-fold moment-kernel, obtaining \(\log\mathfrak{M}^{(q)}(\theta)\) and the scalar residual
  \[
  \mathsf{Anom}^{(q)}_{\mathrm{mom}}(K;\theta)=\log\mathfrak{M}^{(q)}(\theta)-q\log\mathfrak{M}^{(1)}(\theta)
  \]
  (Definition \ref{def:pom-moment-anom-signature}; this is precisely the residual label from the commutation of \(E_m\) with \(\mathrm{MOM}_q\) via outer extraction, Proposition \ref{prop:pom-em-mom-residual-swap}).
  \item \textbf{Rank certificate:} Verify the minimal realization dimension via the Hankel rank lower bound (script \texttt{scripts/exp\_fold\_collision\_moment\_hankel\_rank.py}; Table \ref{tab:fold_collision_moment_hankel_rank}).
  \end{itemize}

The outputs of the above three classes of scripts collectively form the closed loop ``\(r_q\) computable \(\Rightarrow\) tail quantile budget computable'': substituting \(r_q\) into the explicit inequality of Remark \ref{rem:pom-quantile-budget-computable} yields the auditable quota upper bound for the prime-register (Corollary \ref{cor:pom-prime-register-quantile-budget}).

\paragraph{Bridge-Kit interface: second-level regression test suite for strictified commutative diagrams}
The preceding three classes of interfaces respectively provide: constant-level residuals (anomaly signatures), exponential-level slow-mode certificates (worst twisted spectral radius ratios), and tail/prime-register budget certificates (controlled by \(r_q\)).
However, in engineering practice, what is truly needed is a ``\emph{one-click bridge-kit usable as an acceptance criterion}'': any modification to the kernel implementation, truncation strategy, character decomposition, or numerical linear algebra optimization must preserve a set of low-dimensional auditable outputs without drift.
\par
To this end, we organize
\[
\bigl(\mathsf{Anom},\ C_r,\ \delta_{\mathrm{spec}},\ \mathcal{E}_1,\ \eta_{\mathrm{tail}},\ \text{SVD fingerprint}\bigr)
\]
into an independent Bridge-Kit interface (with the option to append the scalar residual of moment outer-extraction commutation \(\mathsf{Anom}^{(q)}_{\mathrm{mom}}\), Definition \ref{def:pom-moment-anom-signature}), outputting a unified JSON second-level summary via script
\texttt{scripts/exp\_bridge\_kit.py}
(\texttt{artifacts/export/bridge\_kit\_real\_input\_40\_arity\_335\_N2.json}).
Here \(C_r\) can be fully aligned with the Dirichlet--Mertens constant tensor of \S\ref{subsec:arity-dirichlet-mertens-tensor}; while \(\delta_{\mathrm{spec}},\mathcal{E}_1,\eta_{\mathrm{tail}}\) and the SVD energy fraction provide an ultra-low-dimensional regression fingerprint (see the \verb|\input|-able snippet generated by the same script):
\input{sections/generated/eq_bridge_kit_real_input_40_arity_335_N2}

\begin{remark}[Strictification of the central extension calibration and constant-level counterterm]\label{rem:pom-bridge-kit-strictification}
In \ref{subsec:pom-anom-cohomology}, this paper has organized the anomaly signature \(\mathsf{Anom}\) as a \(2\)-cohomology class of \(\mathbf{PW}\) (the curvature \(2\)-cell of the fundamental square). Equivalently, it can be viewed as an explicit \(2\)-morphism label on a central extension
\(
\widetilde{\mathbf{PW}}=\mathbf{PW}\ltimes \mathcal{A}_G
\)
(\(\mathcal{A}_G\cong\RR^{\widehat G\setminus\{\chi_0\}}\) being the additive group of ``nontrivial character constant drifts''): thus the ``commutation failure'' of \(E_m\) and \(\mathrm{LIFT}_G\) is no longer an equality violation but rather additive coherence data.
\par
Furthermore, since \(\mathsf{Anom}\) in the calibration of this paper occurs only at the finite-part/constant level (\(\log\mathfrak{M}_\chi-\log\mathfrak{M}_{\chi_0}\)), a pure constant counterterm can strictly cancel this residual at the finite-part interface (anomaly cancellation): it introduces only constant-level padding (e.g., the \(e^{A_\infty}\) in the prime-register budget) without altering the leading exponential slope.
\end{remark}
